\section{Units of Relative Norm One and Planar Point Sets}
\label{geo:section}

This section and Sections~\ref{fw:section} and~\ref{pf:section} turn
arithmetic data of a quadratic extension into planar point sets. Let $K/F$ be a quadratic extension of
number fields, let $\iota$ be the nontrivial automorphism of $K$ over
$F$, and let $(b,c)$ be the signature of $F$. Throughout these three
sections we assume:
\begin{enumerate}
\item[(G1)] $K$ is totally imaginary, and $b\ge1$;
\item[(G2)] $K/F$ is unramified at every finite place;
\item[(G3)] $\rd(K)=\rd(F)=\lambda$.
\end{enumerate}
Here the root discriminant $\rd$ and the discriminants $\Delta_M$ are as in
Section~\ref{tw:tower}. By (G2), $|\Delta_K|=|\Delta_F|^2$, so (G3) amounts
to $\rd(F)=\lambda$.
Put $d=[F:\Q]=b+2c$, $\theta=c/d$ and $\ell=\log\lambda$. Then
$[K:\Q]=2d$, $|\Delta_K|=\lambda^{2d}$, $|\Delta_F|=\lambda^d$ and
$0\le\theta<1/2$, where $\Delta_K$ and $\Delta_F$ are the absolute
discriminants. Theorem~\ref{tw:field-family} provides such extensions,
with $\lambda=2^{9/4}\sqrt{3615}$ and $\theta\ge\theta_*$; there $\iota$
is the complex conjugation $\iota_1$. The arguments of these three
sections use no Galois structure over $B$ or over $\Q$. For a nonzero
fractional ideal $\mathfrak a$ of $K$ or of $F$, $N\mathfrak a$ denotes
its absolute norm, so that $N(x\mathcal O_K)=|N_{K/\Q}x|$
\cite[Propositions~4.1 and~4.2]{MilneANT}. We take the unit theorem and
the covolumes of ideal lattices from \cite{MilneANT}, and the Herbrand
quotient, Hilbert's Theorem~90 and the analytic class-number formula
from \cite{MilneCFT}.

The argument has five steps. Proposition~\ref{geo:mass} is the
class-number formula for the units of relative norm one: it expresses
their regulator, together with the relative class number and the
capitulation kernel, through $L_F(1)$. Lemma~\ref{geo:selection} uses the
split primes and one ideal class to produce many elements
$\beta_1,\dots,\beta_t$ of relative norm one in a single fractional ideal,
with pairwise disjoint cosets $\beta_i\mathcal O_K^1$.
Lemmas~\ref{geo:unfold-lemma} and~\ref{geo:joint-average-lemma} average
over the units of relative norm one and over translates of the ideal
lattice; this counts exactly, on average, the pairs of lattice points in a
window (Definition~\ref{geo:window}) that differ by the image of an element
of $\bigcup_i\beta_i\mathcal O_K^1$. Lemmas~\ref{geo:dual}
and~\ref{geo:poisson} bound the number of lattice points in a window
uniformly, by Poisson summation. Finally, Proposition~\ref{geo:transfer},
the geometric transfer, takes the window to be a sublevel set of the total
energy of a pair of profiles (Definition~\ref{geo:profiles} and
Lemma~\ref{geo:concentration}) and combines these steps.

\subsection{Coordinates}
A real place of $F$ extends to $K$ either as two real places or as one
complex place, and by (G1) the second case occurs. A complex place of
$F$ has two extensions to $K$, since $\C$ has no quadratic extension.
Thus $K$ has exactly $d$ complex places. For each real place $v$ of
$F$, let $\phi_v\colon K\to\C$ be one of the two embeddings inducing
the place of $K$ above $v$. Number the complex places of $F$ from $1$
to $c$. For the $j$-th one fix an embedding $\phi_j^+\colon K\to\C$
whose restriction $\tau_j$ to $F$ induces it, and put
$\phi_j^-=\phi_j^+\circ\iota$. The embeddings $\phi_j^\pm$ induce
the two places of $K$ above this place: they are distinct, and they are
not complex conjugate, because $\tau_j$ is not real. We identify
$K\otimes_\Q\R$ with $\C^d$ through the map $x\mapsto x_\infty$, where
\[
 x_\infty=\bigl((\phi_v(x))_v,\ (\phi_j^+(x),\phi_j^-(x))_{j}\bigr),
\]
write $\mathfrak a_\infty=\{x_\infty:x\in\mathfrak a\}$ for
$\mathfrak a\subset K$, and identify $\C^d$ with $\R^{2d}$ through real
and imaginary parts. Volumes are Lebesgue measure in these coordinates,
and $\langle x,\xi\rangle=\operatorname{Re}\sum_w x_w\overline{\xi_w}$
is the real inner product, where $w$ runs over the $d$ coordinates. For a
Borel set $\Omega$, $|\Omega|$ denotes its volume; for a finite set, $|\cdot|$
denotes its cardinality.

At a real place $v$ of $F$, the automorphism $\iota$ induces complex
conjugation on the completion $\C$ of $K$, so
$\phi_v(\iota x)=\overline{\phi_v(x)}$. Writing $v(y)\in\R$ for the
image of $y\in F$ under the real embedding $v$, we obtain for $x\in K$
\begin{equation}\label{geo:norm-modulus}
 |\phi_v(x)|^2=v(N_{K/F}x),\qquad
 \phi_j^+(x)\,\phi_j^-(x)=\tau_j(N_{K/F}x).
\end{equation}
For $x\in K^\times$ put
\[
 l_j(x)=\tfrac12\bigl(\log|\phi_j^+(x)|-\log|\phi_j^-(x)|\bigr),
 \qquad l(x)=(l_1(x),\dots,l_c(x))\in\R^c .
\]
If $N_{K/F}\beta=1$, then by \eqref{geo:norm-modulus}
$|\phi_v(\beta)|=1$ for every real place $v$, and
$|\phi_j^\pm(\beta)|=e^{\pm l_j(\beta)}$ for every $j$. Accordingly the
$b$ coordinates $\phi_v$ are called the \emph{compact coordinates}:
there the elements of relative norm one lie on the unit circle. The $c$
pairs $(\phi_j^+,\phi_j^-)$ are the \emph{pair coordinates}: there
they have reciprocal moduli. A compact coordinate, or a pair of
coordinates, is called a \emph{block}. The planar sets will be images
under one compact coordinate, which is injective on $K$ and sends every
element of relative norm one to a unit vector. We count ordered pairs of
points at distance one, and divide by two at the end.

\subsection{The Class-Number Formula for Units of Relative Norm One}

Let $\mathcal O_K^\times$ and $\mathcal O_F^\times$ be the unit groups,
and put
\[
 \mathcal O_K^1=\ker\bigl(N_{K/F}\colon\mathcal O_K^\times\to
 \mathcal O_F^\times\bigr),
 \qquad I_N=[\mathcal O_F^\times:N_{K/F}(\mathcal O_K^\times)].
\]
Let $w_K$ be the number of roots of unity in $K$, let $h_K$ and $h_F$ be
the class numbers of $K$ and $F$, put $h_{\rm rel}=h_K/h_F$, let $\kappa$
be the order of the \emph{capitulation kernel}
$\ker\bigl(\Cl(F)\to\Cl(K)\bigr)$, the map being induced by extension of
ideals, and let
$L_F(s)=\zeta_K(s)/\zeta_F(s)$. At $s=1$, $L_F(1)$ denotes the quotient
of the residues of $\zeta_K$ and $\zeta_F$. The regulators
$\operatorname{Reg}_K$ and $\operatorname{Reg}_F$ are the ordinary ones,
which use twice the logarithm of the absolute value at a complex place
\cite[Chapter~5]{MilneANT}.

\begin{proposition}\label{geo:mass}
The roots of unity of $K$ lie in $\mathcal O_K^1$, and the restriction of
$l$ to $\mathcal O_K^1$ has kernel of order $w_K$ and image a full
lattice in $\R^c$. Let $\operatorname{Reg}^1$ be the covolume of this
lattice, with $\operatorname{Reg}^1=1$ if $c=0$. Then
\[
 \kappa=I_N/2^{b-1},\qquad
 \operatorname{Reg}_K/\operatorname{Reg}_F=2^{c-1}I_N\operatorname{Reg}^1,
\]
and consequently
\begin{equation}\label{geo:mass-density}
 \frac{w_K}{h_{\rm rel}\kappa\operatorname{Reg}^1}
 =\frac{2^{d-1}\pi^{b+c}}{\lambda^{d/2}L_F(1)}.
\end{equation}
\end{proposition}

\begin{proof}
\emph{Signs.} Fix a real place $v_0$ of $F$. By
\eqref{geo:norm-modulus}, $v_0(N_{K/F}x)=|\phi_{v_0}(x)|^2>0$ for
$x\in K^\times$, so $-1$ is not a norm from $K^\times$. If $\zeta\in K$
is a root of unity, then $N_{K/F}\zeta$ is a root of unity of $F$,
hence $\pm1$ because $F$ has a real place, and it is positive at $v_0$.
Thus $N_{K/F}\zeta=1$.

\emph{Logarithmic embeddings.} For $x\in\mathcal O_K^\times$ and
$y\in\mathcal O_F^\times$ let
\[
 \operatorname{Log}_K(x)=\bigl((2\log|\phi_vx|)_v,\,(2\log|\phi_j^+x|,
 2\log|\phi_j^-x|)_j\bigr),\qquad
 \operatorname{Log}_F(y)=\bigl((\log|v(y)|)_v,\,(2\log|\tau_j(y)|)_j\bigr).
\]
Their images are lattices in the hyperplanes of coordinate sum zero,
the kernel of $\operatorname{Log}_F$ is $\{\pm1\}$
\cite[Proposition~5.8]{MilneANT}, and $\operatorname{Reg}_K$ and
$\operatorname{Reg}_F$ are the covolumes of the images after one
coordinate is deleted; the choice of the deleted coordinate does not
matter.

\emph{The lattice $l(\mathcal O_K^1)$.} Let $\beta\in\mathcal O_K^1$. Then
$|\phi_v(\beta)|=1$ and $|\phi_j^\pm(\beta)|=e^{\pm l_j(\beta)}$, so
\[
 \operatorname{Log}_K(\beta)=\bigl(0,\ (2l_j(\beta),-2l_j(\beta))_j\bigr).
\]
Hence $l(\beta)=0$ exactly when every conjugate of $\beta$ has absolute
value one, that is, when $\beta$ is a root of unity
\cite[Corollary~5.6]{MilneANT}. By the first step the kernel of $l$ on
$\mathcal O_K^1$ has order $w_K$. The norm of $y\in\mathcal O_F^\times$
is $y^2$, so $N_{K/F}(\mathcal O_K^\times)$ has finite index in
$\mathcal O_F^\times$ and rank $b+c-1$. By the unit theorem
\cite[Theorem~5.1]{MilneANT}, $\mathcal O_K^\times$ has rank $d-1$, so
$\mathcal O_K^1$ has rank $c$. As $\operatorname{Log}_K(\mathcal O_K^1)$
is discrete, so is $l(\mathcal O_K^1)$, and it is a full lattice in
$\R^c$.

\emph{The capitulation kernel.} The Herbrand quotient
$h(M)=|\widehat H^0(\langle\iota\rangle,M)|/|H^1(\langle\iota\rangle,M)|$
\cite[Chapter~II, \S3]{MilneCFT} of the units satisfies
$2\,h(\mathcal O_K^\times)=\prod_v[K_w:F_v]$, where $v$ runs over the
archimedean places of $F$ and $w$ lies above $v$
\cite[Chapter~VII, Proposition~3.1]{MilneCFT}. Each real place
contributes $2$ and each complex place $1$, so
$h(\mathcal O_K^\times)=2^{b-1}$. Since
$\widehat H^0(\langle\iota\rangle,\mathcal O_K^\times)=\mathcal
O_F^\times/N_{K/F}(\mathcal O_K^\times)$ has order $I_N$, we get
$|H^1(\langle\iota\rangle,\mathcal O_K^\times)|=I_N/2^{b-1}$. This group
is $\mathcal O_K^1$ modulo the units $\gamma/\iota\gamma$ with
$\gamma\in\mathcal O_K^\times$. If $\mathfrak a$ is a fractional ideal
of $F$ with $\mathfrak a\mathcal O_K=x\mathcal O_K$, the class of the
unit $x/\iota x$ depends neither on the choice of $x$ nor on the choice
of $\mathfrak a$ in its ideal class. This defines a homomorphism from
$\ker(\Cl(F)\to\Cl(K))$ to $H^1(\langle\iota\rangle,\mathcal O_K^\times)$.
It is injective: if $x/\iota x=\gamma/\iota\gamma$, then $y=x/\gamma$
lies in $F$ and $\mathfrak a\mathcal O_K=y\mathcal O_K$, so
$\mathfrak a=y\mathcal O_F$ by unique factorization. It is surjective:
by Hilbert's Theorem~90 \cite[Chapter~II, Corollary~1.23]{MilneCFT}, an
element of $\mathcal O_K^1$ has the form $x/\iota x$ with
$x\in K^\times$, and then the fractional ideal
$x\mathcal O_K=(\iota x)\mathcal O_K$ is invariant under $\iota$. By
(G2), every $\iota$-invariant fractional ideal of $K$ is extended from
$F$: at a split prime $\mathfrak P\,\iota\mathfrak P$ invariance forces
equal exponents, an inert prime of $F$ stays prime, and no prime
ramifies. Hence $\kappa=I_N/2^{b-1}$.

\emph{Regulators.} Delete the coordinate of $v_0$ from both logarithmic
embeddings, and let $\mathcal U_K\subset\R^{d-1}$ and
$\mathcal U_F\subset\R^{b+c-1}$ be the resulting lattices, of covolumes
$\operatorname{Reg}_K$ and $\operatorname{Reg}_F$. By
\eqref{geo:norm-modulus}, $\operatorname{Log}_F(N_{K/F}x)$ is obtained
from $\operatorname{Log}_K(x)$ by keeping the coordinates of the real
places and adding the two coordinates of each pair. This induces a
linear map $T\colon\R^{d-1}\to\R^{b+c-1}$ such that
$T(\mathcal U_K)$ is the image of $N_{K/F}(\mathcal O_K^\times)$ in
$\mathcal U_F$. The substitution $(x_j,y_j)\mapsto(x_j,x_j+y_j)$ on each
pair has determinant one and turns $T$ into a coordinate projection
whose kernel consists of the $c$ coordinates $x_j$. The covolume of a
lattice in a product of two coordinate spaces is the covolume of its
intersection with the first factor times the covolume of its projection
to the second, whenever both are lattices.

A unit $x\in\mathcal O_K^\times$ lies in the kernel of $T$ exactly
when $\operatorname{Log}_F(N_{K/F}x)=0$, because the deleted coordinate
is minus the sum of the others. Then $N_{K/F}x=\pm1$, and $N_{K/F}x=1$
by the first step. So the intersection of $\mathcal U_K$ with the kernel is
the image of $\mathcal O_K^1$, which after the substitution is
$2l(\mathcal O_K^1)$, of covolume $2^c\operatorname{Reg}^1$. Since
$-1\notin N_{K/F}(\mathcal O_K^\times)$, the projection
$T(\mathcal U_K)$ has index
$[\mathcal O_F^\times:\{\pm1\}N_{K/F}(\mathcal O_K^\times)]=I_N/2$ in
$\mathcal U_F$. Hence
$\operatorname{Reg}_K=2^c\operatorname{Reg}^1\cdot\operatorname{Reg}_FI_N/2$.

\emph{The class-number formula.} The analytic class-number formula
\cite[Chapter~V, Theorem~2.4]{MilneCFT} (compare
\eqref{an:class-formula}), applied to $K$ (with $r_1=0$
and $r_2=d$) and to $F$ (with $r_1=b$, $r_2=c$ and $w_F=2$), gives
\[
 L_F(1)=\frac{(2\pi)^dh_K\operatorname{Reg}_K}{w_K\lambda^d}\cdot
 \frac{2\lambda^{d/2}}{2^b(2\pi)^ch_F\operatorname{Reg}_F}
 =\frac{2^{c+1}\pi^{b+c}h_{\rm rel}}{w_K\lambda^{d/2}}\,
 \frac{\operatorname{Reg}_K}{\operatorname{Reg}_F}.
\]
Substituting
$\operatorname{Reg}_K/\operatorname{Reg}_F=2^{c-1}I_N\operatorname{Reg}^1
=2^{b+c-2}\kappa\operatorname{Reg}^1$ gives \eqref{geo:mass-density}.
\end{proof}

\subsection{Norm-One Elements from Split Primes}

\begin{definition}\label{geo:uniform-types}
Let $\mathcal R$ be a finite set of rational primes. The extension $K/F$
has \emph{uniform local types} $(e_r,f_r)_{r\in\mathcal R}$ if, for every
$r\in\mathcal R$, every prime of $F$ above $r$ splits in $K$ and has
absolute type $(e_r,f_r)$ (Definition~\ref{tw:type}).
\end{definition}

If $K/F$ has uniform local types $(e_r,f_r)_{r\in\mathcal R}$, then $F$
has exactly $d/(e_rf_r)$ primes above $r$, each of norm $r^{f_r}$, since
the products $e_{\mathfrak p}f_{\mathfrak p}$ of the absolute ramification
indices and residue degrees of the primes $\mathfrak p$ of $F$ above $r$
add up to $[F:\Q]=d$. Given integers $k_r\ge0$, define
\begin{equation}\label{geo:HJ}
 H=\sum_{r\in\mathcal R}\frac{k_r\log r}{e_r},\qquad
 J=\sum_{r\in\mathcal R}\frac{\log(k_r+1)}{e_rf_r}.
\end{equation}

\begin{lemma}\label{geo:selection}
Assume that $K/F$ has uniform local types $(e_r,f_r)_{r\in\mathcal R}$, and
fix integers $k_r\ge0$. There are a
fractional ideal $I$ of $K$ with $N(I)=e^{-dH}$, an integer
$t\ge e^{dJ}/(\kappa h_{\rm rel})$, and elements
$\beta_1,\dots,\beta_t\in I$ of relative norm one whose cosets
$\beta_i\mathcal O_K^1$ are pairwise disjoint.
\end{lemma}

\begin{proof}
For each prime $\mathfrak p$ of $F$ above a prime $r\in\mathcal R$, write
$\mathfrak p\mathcal O_K=\mathfrak P_{\mathfrak p}\,\iota\mathfrak
P_{\mathfrak p}$ and $k_{\mathfrak p}=k_r$. For integer vectors
$\boldsymbol j=(j_{\mathfrak p})$ with $0\le j_{\mathfrak p}\le
k_{\mathfrak p}$ put
\[
 \mathfrak A_{\boldsymbol j}=\prod_{\mathfrak p}
 \mathfrak P_{\mathfrak p}^{j_{\mathfrak p}},\qquad
 \mathfrak J_{\boldsymbol j}=\prod_{\mathfrak p}
 \mathfrak P_{\mathfrak p}^{j_{\mathfrak p}}
 (\iota\mathfrak P_{\mathfrak p})^{k_{\mathfrak p}-j_{\mathfrak p}}.
\]
There are $\prod_{\mathfrak p}(k_{\mathfrak p}+1)=e^{dJ}$ such vectors.
The quotient of $\Cl(K)$ by the image of $\Cl(F)$ has order
$h_K\kappa/h_F=\kappa h_{\rm rel}$, so one class of this quotient
contains the classes of $\mathfrak A_{\boldsymbol j}$ for a set of
$t\ge e^{dJ}/(\kappa h_{\rm rel})$ vectors $\boldsymbol j$. Fix one of
them, $\boldsymbol j_0$, and put $I=\mathfrak J_{\boldsymbol j_0}^{-1}$.
For each $\boldsymbol j$ in the set, write
$\mathfrak A_{\boldsymbol j}\mathfrak A_{\boldsymbol j_0}^{-1}
=x\,\mathfrak a\mathcal O_K$ with $x\in K^\times$ and $\mathfrak a$ a
fractional ideal of $F$, and put $\beta_{\boldsymbol j}=x/\iota x$.
Then $N_{K/F}\beta_{\boldsymbol j}=1$, and since $\iota$ fixes
$\mathfrak a\mathcal O_K$,
\[
 \beta_{\boldsymbol j}\mathcal O_K
 =\mathfrak A_{\boldsymbol j}\mathfrak A_{\boldsymbol j_0}^{-1}
 \bigl(\iota\mathfrak A_{\boldsymbol j}\bigr)^{-1}
 \iota\mathfrak A_{\boldsymbol j_0}
 =\mathfrak J_{\boldsymbol j}I.
\]
As $\mathfrak J_{\boldsymbol j}$ is integral,
$\beta_{\boldsymbol j}\in I$. The ideals $\mathfrak J_{\boldsymbol j}I$
are distinct for distinct $\boldsymbol j$, and every element of
$\beta_{\boldsymbol j}\mathcal O_K^1$ generates
$\mathfrak J_{\boldsymbol j}I$, so the cosets are disjoint. Finally,
since each $\mathfrak p$ splits in $K$,
$N(I)^{-1}=\prod_{\mathfrak p}(N\mathfrak p)^{k_{\mathfrak p}}
=\prod_r r^{f_rk_rd/(e_rf_r)}=e^{dH}$.
\end{proof}

\subsection{Unit Averaging and Translation}
In this subsection and the next, $I$ is a fractional ideal of $K$ and
$\beta_1,\dots,\beta_t\in I$ are elements of relative norm one whose
cosets $\beta_i\mathcal O_K^1$ are pairwise disjoint;
Lemma~\ref{geo:selection} provides such data. For $h\in\R^c$ let
$W_h\colon\C^d\to\C^d$ multiply the coordinate $\phi_j^+$ by
$e^{-h_j}$ and $\phi_j^-$ by $e^{h_j}$, and fix the compact
coordinates. It is a real diagonal map of determinant one. Put
$\Lambda_h=W_hI_\infty$. For a lattice $\Lambda\subset\C^d$,
$\covol(\Lambda)$ denotes the volume of a fundamental domain. By
\cite[Proposition~4.26]{MilneANT}, extended to fractional ideals by
scaling and \cite[Proposition~4.2]{MilneANT},
\begin{equation}\label{geo:covolume}
 \covol(\mathfrak a_\infty)=2^{-d}|\Delta_K|^{1/2}N\mathfrak a
\end{equation}
for every fractional ideal $\mathfrak a$ of $K$. Hence
$\covol(\Lambda_h)=2^{-d}\lambda^dN(I)$, independently of $h$; for the
ideal of Lemma~\ref{geo:selection},
\begin{equation}\label{geo:determinant}
 \covol(\Lambda_h)=2^{-d}\lambda^de^{-dH}.
\end{equation}
If $N_{K/F}\beta=1$, then $W_h\beta_\infty$ has compact coordinates of
modulus one and pair coordinates of moduli $e^{\pm(l_j(\beta)-h_j)}$.
Let $\Pi^1$ be a fundamental parallelepiped of the lattice
$l(\mathcal O_K^1)$; its volume is $\operatorname{Reg}^1$. When $c=0$,
integrals over $\R^0$ and over $\Pi^1$ are evaluations at the single
point.

\begin{lemma}\label{geo:unfold-lemma}
For every Borel function $\Phi\colon\R^c\to[0,\infty]$,
\begin{equation}\label{geo:unfold}
 \frac1{\operatorname{Reg}^1}\int_{\Pi^1}\sum_{i=1}^t
 \sum_{\beta\in\beta_i\mathcal O_K^1}\Phi(l(\beta)-h)\,dh
 =\frac{tw_K}{\operatorname{Reg}^1}\int_{\R^c}\Phi(u)\,du .
\end{equation}
\end{lemma}

\begin{proof}
The map $l$ is a homomorphism on $K^\times$, so
$l(\beta_i\mathcal O_K^1)=l(\beta_i)+l(\mathcal O_K^1)$, and by
Proposition~\ref{geo:mass} each point of this translate of
$l(\mathcal O_K^1)$ is the image of exactly $w_K$ elements of
$\beta_i\mathcal O_K^1$. The translates $\Pi^1-x$, with
$x\in l(\mathcal O_K^1)$, tile $\R^c$, so by Tonelli's theorem
$\int_{\Pi^1}\sum_{x\in l(\mathcal
O_K^1)}\Phi(l(\beta_i)+x-h)\,dh=\int_{\R^c}\Phi(u)\,du$ for each
$i$.
\end{proof}

\begin{definition}\label{geo:window}
A \emph{window} is a bounded Borel set $\Omega\subset\C^d$ of positive
volume that is invariant under rotation of each coordinate separately.
For $u\in\R^c$ let $\nu(u)\in\C^d$ have compact coordinates $1$ and
pair coordinates $(e^{u_j},e^{-u_j})$, and for a Borel set
$\Omega\subset\C^d$ put
\[
 \Phi_\Omega(u)=|\Omega\cap(\Omega-\nu(u))|.
\]
\end{definition}

\begin{lemma}\label{geo:phi-omega}
Let $\Omega$ be a window, let $h\in\R^c$, and let $\beta\in K$ satisfy
$N_{K/F}\beta=1$. Then
$|\Omega\cap(\Omega-W_h\beta_\infty)|=\Phi_\Omega(l(\beta)-h)$.
\end{lemma}

\begin{proof}
If $x\in\C^d$ has compact coordinates of modulus one and pair
coordinates of moduli $e^{\pm u_j}$, then a rotation of the coordinates
preserves $\Omega$ and maps $\nu(u)$ to $x$, so
$|\Omega\cap(\Omega-x)|=\Phi_\Omega(u)$. By the remark before
Lemma~\ref{geo:unfold-lemma}, $x=W_h\beta_\infty$ is such a point, with
$u=l(\beta)-h$.
\end{proof}

For a window $\Omega$, the function $\Phi_\Omega$ is continuous, bounded
by $|\Omega|$, and vanishes when some $e^{|u_j|}$ exceeds the diameter of
$\Omega$: if $x$ and $x+\nu(u)$ both lie in $\Omega$, then $e^{|u_j|}$,
the modulus of a coordinate of their difference $\nu(u)$, is at most the
diameter of $\Omega$.

Let $\Omega$ be a window. Fix a $\Z$-basis of $I_\infty$, and for
$y\in[0,1)^{2d}$ let $y_I\in\C^d$ be the point with coordinate vector $y$
in this basis. For $h\in\R^c$ let
$\mathcal X(h,y)=(\Lambda_h+W_hy_I)\cap\Omega$ and
$n(h,y)=|\mathcal X(h,y)|$, and let $\mathcal E(h,y)$ be the number of
\emph{ordered edges} of $\mathcal X(h,y)$, that is, of pairs $(x,\beta)$
with $x\in\mathcal X(h,y)$, $\beta\in\bigcup_i\beta_i\mathcal O_K^1$ and
$x+W_h\beta_\infty\in\Omega$. Since $\beta\in I$, the point
$x+W_h\beta_\infty$ then lies in $\mathcal X(h,y)$.

\begin{lemma}\label{geo:joint-average-lemma}
Let $\Omega$ be a window and suppose $n(h,y)\le n_{\max}<\infty$ for
all $(h,y)$. Then there is a pair $(h,y)$ with $n(h,y)>0$ and
\begin{equation}\label{geo:joint-average}
 \frac{\mathcal E(h,y)}{n(h,y)}\ge
 \frac{tw_K}{\operatorname{Reg}^1}\,
 \frac{\int_{\R^c}\Phi_\Omega(u)\,du}{|\Omega|}.
\end{equation}
\end{lemma}

\begin{proof}
For fixed $x$, distinct $\beta$ give distinct points
$x+W_h\beta_\infty\ne x$, so $\mathcal E\le n^2\le n_{\max}^2$. Both $n$
and $\mathcal E$ are countable sums of indicator functions of Borel sets
in $(h,y)$. For fixed $h$, the map $y\mapsto W_hy_I$ has Jacobian
$\covol(\Lambda_h)$ and maps $[0,1)^{2d}$ onto a fundamental domain of
$\Lambda_h$, so tiling and Lemma~\ref{geo:phi-omega} give
\[
 \int_{[0,1)^{2d}}n(h,y)\,dy=\frac{|\Omega|}{\covol(\Lambda_h)},\qquad
 \int_{[0,1)^{2d}}\mathcal E(h,y)\,dy
 =\frac1{\covol(\Lambda_h)}\sum_{i}\sum_{\beta\in\beta_i\mathcal O_K^1}
 \Phi_\Omega(l(\beta)-h).
\]
Average over $h\in\Pi^1$ with respect to $dh/\operatorname{Reg}^1$ and
apply \eqref{geo:unfold}. For the probability measure
$dh\,dy/\operatorname{Reg}^1$ on $\Pi^1\times[0,1)^{2d}$ this gives
$\mathbb E\,\mathcal E=\rho\,\mathbb En$ and
$\mathbb En=|\Omega|/\covol(\Lambda_h)>0$, where $\rho$ is the right
side of \eqref{geo:joint-average}. If $\mathcal E<\rho n$ held wherever
$n>0$, then, since $\mathcal E=0$ where $n=0$ and $\{n>0\}$ has positive
measure, we would get $\mathbb E\,\mathcal E<\rho\,\mathbb En$.
\end{proof}

\begin{lemma}\label{geo:planar}
Let $\Omega$ be a window, let $(h,y)$ be a pair, and let $v_0$ be a real
place of $F$. For $x\in\C^d$ let $x_{v_0}$ be its compact coordinate at
$v_0$, and let $U=\{x_{v_0}:x\in\mathcal X(h,y)\}\subset\C=\R^2$. Then
$|U|=n(h,y)$ and $u(U)\ge\mathcal E(h,y)/2$.
\end{lemma}

\begin{proof}
Two distinct points of $\mathcal X(h,y)$ differ by $W_h\gamma_\infty$
with $0\ne\gamma\in I$, whose $v_0$-coordinate $\phi_{v_0}(\gamma)$ is
nonzero, so $|U|=n(h,y)$. An ordered edge $(x,\beta)$ gives two points of
$U$ at distance $|\phi_{v_0}(\beta)|=1$, and distinct ordered edges give
distinct ordered pairs of points. Hence $u(U)\ge\mathcal E(h,y)/2$.
\end{proof}

\subsection{Uniform Lattice Counting}
For a lattice $\Lambda\subset\C^d$ let
$\Lambda^*=\{\xi\in\C^d:\langle\xi,x\rangle\in\Z\text{ for all
}x\in\Lambda\}$ be its dual lattice, and let
$\|\xi\|_*=\sum_w|\xi_w|$, a norm on $\C^d$. We use the Fourier
transform $\widehat\Psi(\xi)=\int_{\C^d}\Psi(x)e^{-2\pi i\langle
x,\xi\rangle}\,dx$.

\begin{lemma}\label{geo:dual}
Let $\mathfrak a$ be a fractional ideal of $K$ and $h\in\R^c$. Every
nonzero $\xi$ in the dual lattice of $W_h\mathfrak a_\infty$ satisfies
\[
 \prod_w|\xi_w|\ge\frac{2^d}{\lambda^d(N\mathfrak a)^{1/2}},
 \qquad
 \|\xi\|_*\ge\frac{2d}{\lambda(N\mathfrak a)^{1/(2d)}}.
\]
In particular, if $N(I)=e^{-dH}$, as for the ideal of
Lemma~\ref{geo:selection}, every nonzero dual vector of $\Lambda_h$
satisfies
\begin{equation}\label{geo:dual-product}
 \prod_w|\xi_w|\ge\mu^d,\qquad \mu=2e^{H/2-\ell},\qquad
 \|\xi\|_*\ge d\mu.
\end{equation}
\end{lemma}

\begin{proof}
First let $h=0$. For $x,y\in K$ we have
$\langle x_\infty,\overline{y_\infty}\rangle
=\operatorname{Re}\sum_w\phi_w(xy)=\tfrac12\Tr_{K/\Q}(xy)$, where the
bar is coordinatewise complex conjugation and $\phi_w$ is the
embedding of the coordinate $w$, because these embeddings and their
complex conjugates are all the embeddings of $K$. Let
$\mathfrak a^\vee=\{y\in K:\Tr_{K/\Q}(y\mathfrak a)\subset\Z\}$. It is an
$\mathcal O_K$-module, and since the trace form is nondegenerate, the
dual basis of a $\Z$-basis of $\mathfrak a$ is a $\Z$-basis of
$\mathfrak a^\vee$. So $\mathfrak a^\vee$ is a fractional ideal, and
the dual lattice of $\mathfrak a_\infty$ is
$\overline{(2\mathfrak a^\vee)_\infty}$. Dual lattices have reciprocal
covolumes, and complex conjugation preserves volume. With
\eqref{geo:covolume} and $N(2\mathfrak a^\vee)=2^{2d}N(\mathfrak a^\vee)$, this
gives $N(\mathfrak a^\vee)=\lambda^{-2d}(N\mathfrak a)^{-1}$. If
$0\ne y\in\mathfrak a^\vee$, then $y\mathcal O_K\subset\mathfrak a^\vee$
and $|N_{K/\Q}y|=N(y\mathcal O_K)\ge N(\mathfrak a^\vee)$
\cite[Propositions~4.1 and~4.2]{MilneANT}. For
$\xi=\overline{(2y)_\infty}$,
\[
 \prod_w|\xi_w|=|N_{K/\Q}(2y)|^{1/2}
 \ge\bigl(2^{2d}\lambda^{-2d}(N\mathfrak a)^{-1}\bigr)^{1/2}.
\]
For general $h$, the dual lattice of $W_h\mathfrak a_\infty$ is the
image of the dual lattice of $\mathfrak a_\infty$ under $W_{-h}$,
because $W_h$ is real diagonal, and $W_{-h}$ preserves
$\prod_w|\xi_w|$. The bound for $\|\xi\|_*$ follows from the
arithmetic-geometric mean inequality. For \eqref{geo:dual-product} take
$\mathfrak a=I$.
\end{proof}

\begin{lemma}\label{geo:poisson}
Let $\Lambda\subset\C^d$ be a lattice such that
$\prod_w|\xi_w|\ge\mu^d$ for every nonzero $\xi\in\Lambda^*$, where
$\mu>0$. Let $\Psi\colon\C^d\to[0,\infty)$ be smooth, bounded and
integrable, with $\int\Psi>0$ and all partial derivatives bounded, and
suppose that
\[
 |\widehat\Psi(\xi)|\le M^de^{-\sigma\|\xi\|_*}\int\Psi
 \qquad(\xi\in\C^d)
\]
for some $M\ge1$ and $\sigma>0$. If the \emph{Fourier condition}
\begin{equation}\label{geo:poisson-condition}
 \sigma\mu\ge\log M+2\log5+2
\end{equation}
holds, then for every $y\in\C^d$,
\begin{equation}\label{geo:pointwise-poisson}
 \sum_{x\in\Lambda+y}\Psi(x)\le\frac2{\covol(\Lambda)}\int\Psi.
\end{equation}
\end{lemma}

\begin{proof}
A nonzero $\xi\in\Lambda^*$ has no zero coordinate, and the
arithmetic-geometric mean inequality gives $\|\xi\|_*\ge d\mu$. Hence
distinct points of $\Lambda^*$ are at $\|\cdot\|_*$-distance at least
$d\mu$, and the open $\|\cdot\|_*$-balls of radius $d\mu/2$ about them
are disjoint. For $j\ge1$, the balls about the points with
$jd\mu\le\|\xi\|_*<(j+1)d\mu$ lie in the ball of radius $(j+3/2)d\mu$
about $0$. Comparing volumes in $\R^{2d}$ bounds the number of these
points by $(2j+3)^{2d}\le5^{2dj}$. By the Fourier condition
\eqref{geo:poisson-condition} and
$M\ge1$,
\[
 M^d\sum_{\xi\in\Lambda^*\setminus\{0\}}e^{-\sigma\|\xi\|_*}
 \le\sum_{j\ge1}e^{d(\log M+2j\log5-j\sigma\mu)}
 \le\sum_{j\ge1}e^{-2dj}\le1 .
\]

For $\varepsilon>0$ put $\Psi_\varepsilon(x)=\Psi(x)e^{-\varepsilon|x|^2}$, a Schwartz
function. Its Fourier transform is $\widehat\Psi*G_\varepsilon$, where
$G_\varepsilon(\zeta)=(\pi/\varepsilon)^de^{-\pi^2|\zeta|^2/\varepsilon}$ is a probability
density on $\R^{2d}$. Hence
$|\widehat{\Psi_\varepsilon}(\xi)|\le m_\varepsilon M^de^{-\sigma\|\xi\|_*}\int\Psi$,
with $m_\varepsilon=\int e^{\sigma\|\zeta\|_*}G_\varepsilon(\zeta)\,d\zeta$. The function
$y\mapsto\sum_{x\in\Lambda}\Psi_\varepsilon(x+y)$ is smooth and
$\Lambda$-periodic, its Fourier coefficient at $\xi\in\Lambda^*$ is
$\widehat{\Psi_\varepsilon}(\xi)/\covol(\Lambda)$, and these coefficients are
absolutely summable. This gives the Poisson summation formula
\[
 \sum_{x\in\Lambda}\Psi_\varepsilon(x+y)
 =\frac1{\covol(\Lambda)}\sum_{\xi\in\Lambda^*}\widehat{\Psi_\varepsilon}(\xi)
 e^{2\pi i\langle\xi,y\rangle}.
\]
Using $\widehat{\Psi_\varepsilon}(0)=\int\Psi_\varepsilon$ and the bound above, the sum
is at most $(\int\Psi_\varepsilon+m_\varepsilon\int\Psi)/\covol(\Lambda)$. As
$\varepsilon\to0$, we have $m_\varepsilon\to1$ by dominated convergence, since
$G_\varepsilon$ is the image of $G_1$ under $\zeta\mapsto\sqrt\varepsilon\,\zeta$;
moreover $\Psi_\varepsilon$ increases to $\Psi$. Monotone convergence gives
\eqref{geo:pointwise-poisson}.
\end{proof}

The next bound is used in Section~\ref{pf:section} for the Gaussian
real-place profile.

\begin{lemma}\label{geo:gaussian}
If $0<a_\R\le1$ and $\Psi_\R(z)=e^{-a_\R|z|^2}$ on $\C$, then
$\widehat{\Psi_\R}(\xi)=(\pi/a_\R)e^{-\pi^2|\xi|^2/a_\R}$ and
$|\widehat{\Psi_\R}(\xi)|\le2e^{-|\xi|}\int\Psi_\R$.
\end{lemma}

\begin{proof}
The formula is the Gaussian integral, and $\int\Psi_\R=\pi/a_\R$. Since
$a_\R\le1$, it suffices that $\pi^2x^2-x+\log2\ge0$ for $x\ge0$, and the
minimum of the left side is $\log2-1/(4\pi^2)>0$.
\end{proof}

\subsection{Profiles at Real and Complex Places}

\begin{definition}\label{geo:profiles}
Fix $0<\delta<1$ and put $p=2/(1+\delta)$, the \emph{Lebesgue exponent}
(not a prime), so that $1<p<2$ and $p(1+\delta)=2$. A \emph{real-place
profile} is a positive radial function $f_\R$ on $\C$, used at the compact
coordinates, and a \emph{complex-place profile} is a positive function
$g_\C$ on $\C^2$ that depends only on $|z|$ and $|w|$, used at the pair
coordinates. Their \emph{masses} $A_\R$, $A_\C$, \emph{overlaps} $I_\R$,
$I_\C$ and \emph{functionals} $J_\R$, $J_\C$ are
\[
 A_\R=\int_\C f_\R^p,\qquad I_\R=\int_\C f_\R(z)f_\R(z+1)\,dz,
 \qquad A_\C=\int_{\C^2}g_\C^p,
\]
\[
 I_\C=\int_\R\int_{\C^2}g_\C(z,w)\,g_\C(z+e^u,w+e^{-u})\,dz\,dw\,du,
\]
\[
 J_\R=\log I_\R-(1+\delta)\log A_\R,\qquad
 J_\C=\log I_\C-(1+\delta)\log A_\C.
\]
When these integrals are finite and positive, the \emph{endpoint laws}
are the probability densities $f_\R(z)f_\R(z+1)/I_\R$ on $\C$ and
$g_\C(z,w)g_\C(z+e^u,w+e^{-u})/I_\C$ on $\C^2\times\R$. The functions
$V_\R=-\log f_\R$ and $V_\C=-\log g_\C$ are the \emph{energies} of the
profiles.
\end{definition}

The exponent $p$ is chosen so that the terms in the mean energy cancel at
the end of the proof of Proposition~\ref{geo:transfer}.

\begin{definition}\label{geo:admissible}
Let $M\ge1$ and $\sigma>0$. The pair $(f_\R,g_\C)$ is \emph{admissible
with Fourier constants $(M,\sigma)$} if:
\begin{enumerate}
\item[(P1)] $f_\R^p$ and $g_\C^p$ are smooth and bounded, and all their
partial derivatives are bounded;
\item[(P2)] $A_\R$, $I_\R$, $A_\C$ and $I_\C$ are finite and positive;
\item[(P3)] $V_\R$ and $V_\C$ are bounded below and coercive, that is,
their sublevel sets are bounded;
\item[(P4)] $V_\R(z)$ and $V_\C(z,w)$ have finite second moments under
the endpoint laws;
\item[(P5)] $|\widehat{f_\R^p}(\xi)|\le Me^{-\sigma|\xi|}A_\R$ and
$|\widehat{g_\C^p}(\xi_1,\xi_2)|\le M^2e^{-\sigma(|\xi_1|+|\xi_2|)}A_\C$
for all $\xi,\xi_1,\xi_2\in\C$.
\end{enumerate}
\end{definition}

Section~\ref{pf:section} proves admissibility for the profiles we use.

\begin{lemma}\label{geo:tensor}
Let $(f_\R,g_\C)$ be admissible with Fourier constants $(M,\sigma)$, and
let $\Psi=\prod_vf_\R^p\prod_jg_\C^p$ be the product function on $\C^d$,
with one factor for each block. Then $\Psi$ satisfies the hypotheses that
Lemma~\ref{geo:poisson} places on $\Psi$, other than the Fourier condition
\eqref{geo:poisson-condition}: it is smooth and bounded, with bounded
derivatives, $\int\Psi=A_\R^bA_\C^c$, and
\[
 |\widehat\Psi(\xi)|\le M^de^{-\sigma\|\xi\|_*}\int\Psi\qquad(\xi\in\C^d).
\]
\end{lemma}

\begin{proof}
The first properties follow from (P1) and (P2). The Fourier transform of
$\Psi$ is the product of the transforms of the factors, so (P5) gives
$|\widehat\Psi(\xi)|\le M^{b+2c}e^{-\sigma\|\xi\|_*}\int\Psi
=M^de^{-\sigma\|\xi\|_*}\int\Psi$.
\end{proof}

\begin{lemma}\label{geo:concentration}
Let $f_\R$ and $g_\C$ be profiles satisfying (P1), (P2) and (P4), and let
$\eta>0$. Let $m_\R$ and $m_\C$ be the means of $V_\R(z)$ and $V_\C(z,w)$
under the endpoint laws, and let $\mathrm{Var}_\R$ and $\mathrm{Var}_\C$ be
their variances. For $x\in\C^d$ with compact coordinates $x_v$ and pair
coordinates $(x_j^+,x_j^-)$, let
$V(x)=\sum_vV_\R(x_v)+\sum_jV_\C(x_j^+,x_j^-)$ be the total energy, put
$\bar m=(bm_\R+cm_\C)/d$, and let
\[
 \Omega=\bigl\{x\in\C^d:V(x)\le d(\bar m+\eta)\bigr\}.
\]
Then, for all $d$ larger than a bound that depends only on $\eta$,
$\mathrm{Var}_\R$ and $\mathrm{Var}_\C$,
\[
 \int_{\R^c}\Phi_\Omega(u)\,du\ge\tfrac12I_\R^bI_\C^ce^{2d(\bar m-\eta)}.
\]
\end{lemma}

\begin{proof}
By (P1), $f_\R$ and $g_\C$ are continuous, so $V$ is continuous and
$\Omega$ is closed, hence a Borel set. The function
$(x,u)\mapsto e^{-V(x)}e^{-V(x+\nu(u))}/(I_\R^bI_\C^c)$ is the
product of the endpoint densities of the blocks, hence a probability
density on $\C^d\times\R^c$. Under it, the energies $V(x)$ and
$V(x+\nu(u))$ of the two endpoints are sums of independent terms, one for
each block, and $V(x)$ has mean $d\bar m$. On a compact block the map
$z\mapsto-z-1$, and on a pair block the map
$(z,w,u_j)\mapsto(-z-e^{u_j},-w-e^{-u_j},u_j)$, preserves the endpoint
density and carries the first endpoint to minus the second. Since the
profiles are even, $V(x+\nu(u))$ has the same law as $V(x)$. By
Chebyshev's inequality, each of the two energies differs from $d\bar m$
by more than $\eta d$ with probability at most
$\max(\mathrm{Var}_\R,\mathrm{Var}_\C)/(\eta^2d)$. So the set
$\mathcal Y\subset\C^d\times\R^c$ where both energies lie in
$[d(\bar m-\eta),d(\bar m+\eta)]$ has probability at least $1/2$ for all
$d$ larger than a bound that depends only on $\eta$, $\mathrm{Var}_\R$ and
$\mathrm{Var}_\C$. On $\mathcal Y$ both endpoints lie in $\Omega$ and
$e^{-V(x)}e^{-V(x+\nu(u))}\le e^{-2d(\bar m-\eta)}$. Therefore, with
$|\mathcal Y|$ the Lebesgue measure of $\mathcal Y$,
\[
 \int_{\R^c}\Phi_\Omega(u)\,du\ge|\mathcal Y|
 \ge e^{2d(\bar m-\eta)}I_\R^bI_\C^c\,\mathbb P(\mathcal Y)
 \ge\tfrac12I_\R^bI_\C^ce^{2d(\bar m-\eta)}.
\]
\end{proof}

\begin{proposition}[Geometric transfer]\label{geo:transfer}
Let $0<\delta<1$ and $C\in\R$. Let $(K_i/F_i)_{i\ge1}$ be quadratic
extensions satisfying (G1)--(G3) with a common $\lambda$ and with
$d_i=[F_i:\Q]\to\infty$, such that $\log L_{F_i}(1)\le d_iC$ and such
that every $K_i/F_i$ has the same uniform local types
$(e_r,f_r)_{r\in\mathcal R}$. Fix integers $k_r\ge0$, let $H$ and $J$ be
as in \eqref{geo:HJ}, and let $\mu=2e^{H/2-\ell}$. Let $(f_\R,g_\C)$ be
admissible with Fourier constants $(M,\sigma)$ satisfying the Fourier
condition \eqref{geo:poisson-condition}. Put
\begin{equation}\label{geo:general-margin}
 \mathcal M(\theta)=J-\delta H-(1/2-\delta)\ell-C
 +(1-\delta)\log2+(1-\theta)\log\pi
 +(1-2\theta)J_\R+\theta J_\C.
\end{equation}
We call $\mathcal M(\theta)$ the \emph{margin}. If
$\inf_i\mathcal M(\theta_i)>0$, where $\theta_i$ is the value of
$\theta$ for $F_i$, then there are finite sets $U_i\subset\R^2$ with
$|U_i|\to\infty$ and $u(U_i)/|U_i|^{1+\delta}\to\infty$.
\end{proposition}

\begin{proof}
\emph{The window.} Let $m_\R$, $m_\C$, $\mathrm{Var}_\R$ and
$\mathrm{Var}_\C$ be as in Lemma~\ref{geo:concentration}; the variances
are finite by (P4). Fix $\eta>0$ with
$4\eta<\inf_i\mathcal M(\theta_i)$, and consider one extension $K/F$ of
the sequence, with $d$ large. Take $I$, $t$ and
$\beta_1,\dots,\beta_t$ from Lemma~\ref{geo:selection}, so that
$\Lambda_h$ is defined; the counts $n(h,y)$ and $\mathcal E(h,y)$ are those
of the window $\Omega$ chosen next. With the total energy $V$ and the mean
$\bar m$ of Lemma~\ref{geo:concentration}, let
\[
 \Omega=\bigl\{x\in\C^d:V(x)\le d(\bar m+\eta)\bigr\}.
\]
By (P1) and (P3) (see the proof of Lemma~\ref{geo:concentration}), this
is a bounded Borel set, and it is invariant under
coordinate rotations because the profiles are radial in each
coordinate. By Lemma~\ref{geo:concentration},
\[
 \int_{\R^c}\Phi_\Omega(u)\,du\ge\tfrac12I_\R^bI_\C^ce^{2d(\bar m-\eta)}
\]
for all large $d$, uniformly in the signature. In particular
$|\Omega|>0$ for large $d$, so $\Omega$ is a window.

\emph{Counting points.} With the product function
$\Psi=\prod_vf_\R^p\prod_jg_\C^p=e^{-pV}$ of Lemma~\ref{geo:tensor}, we
have $\mathbf 1_\Omega\le e^{pd(\bar m+\eta)}\Psi$, so
\[
 |\Omega|\le e^{pd(\bar m+\eta)}A_\R^bA_\C^c,
\]
and Lemmas~\ref{geo:dual}, \ref{geo:poisson} and~\ref{geo:tensor} give,
for all $(h,y)$,
\[
 n(h,y)\le\frac{2e^{pd(\bar m+\eta)}A_\R^bA_\C^c}{\covol(\Lambda_h)}
 =2e^{dY_\eta},\qquad
 Y_\eta=H+\log2-\ell+(1-2\theta)\log A_\R+\theta\log A_\C
 +p(\bar m+\eta),
\]
by \eqref{geo:determinant}, $b/d=1-2\theta$ and $c/d=\theta$.

\emph{Counting edges.} Now take a pair $(h,y)$ from
Lemma~\ref{geo:joint-average-lemma}, applied with
$n_{\max}=2e^{dY_\eta}$. By $t\ge e^{dJ}/(\kappa h_{\rm rel})$,
Proposition~\ref{geo:mass}, $\log L_F(1)\le dC$ and $b+c=(1-\theta)d$,
\[
 \frac{tw_K}{\operatorname{Reg}^1}\ge
 \frac{e^{dJ}2^{d-1}\pi^{b+c}}{\lambda^{d/2}L_F(1)}
 \ge\tfrac12\exp\bigl(d[J+\log2+(1-\theta)\log\pi-\ell/2-C]\bigr).
\]
Combined with the bounds for $|\Omega|$ and for the integral of
$\Phi_\Omega$, \eqref{geo:joint-average} shows that the set
$\mathcal X=\mathcal X(h,y)$ has $n=|\mathcal X|>0$ points and at least
$\tfrac14e^{dX_\eta}n$ ordered edges, where
\[
 \begin{aligned}
 X_\eta={}&J+\log2+(1-\theta)\log\pi-\ell/2-C
 +(1-2\theta)(\log I_\R-\log A_\R)\\
 &+\theta(\log I_\C-\log A_\C)+(2-p)\bar m-(2+p)\eta.
 \end{aligned}
\]
Since $n\le2e^{dY_\eta}$ and $p(1+\delta)=2$, the terms in $\bar m$ cancel,
and
\[
 \frac{\mathcal E(h,y)}{n^{1+\delta}}
 \ge\frac{e^{dX_\eta}}{4\cdot2^\delta e^{\delta dY_\eta}}
 =2^{-2-\delta}e^{d(\mathcal M(\theta)-4\eta)}.
\]

\emph{The planar set.} Finally let $v_0$ be a real place of $F$, and let
$U=\{x_{v_0}:x\in\mathcal X\}\subset\C=\R^2$. By
Lemma~\ref{geo:planar}, $|U|=n$ and $u(U)\ge\mathcal E(h,y)/2$, so
\[
 \frac{u(U)}{|U|^{1+\delta}}
 \ge2^{-3-\delta}e^{d(\mathcal M(\theta)-4\eta)},
\]
which tends to infinity along the sequence. Since $u(U)\le|U|^2/2$ and
$\delta<1$, also $|U|\to\infty$. These sets are the $U_i$. The
conclusion concerns a sequence of cardinalities, not every large
cardinality.
\end{proof}
